\documentclass{article}
\usepackage[margin=1.5cm]{geometry}
\usepackage{amsmath}
\usepackage{enumitem}

\setlist{nosep,leftmargin=*}

\begin{document}
\twocolumn

\title{Chapters 6--7 Warm-Up: Functions and Modules}
\author{Prof. Jordan C. Hanson}
\date{September 21, 2026}

\maketitle

\section{Starting Data}

\begin{enumerate}

\item We will review functions and modules by extending the
patient-database ideas from Problem Set 4.  Suppose a patient record is represented by a list:

\begin{verbatim}
patient = [1042,"Maria Garcia",37]
\end{verbatim}

The entries represent the patient ID, name, and age.

\end{enumerate}

\section{Functions and Parameters}

\begin{enumerate}

\item Write a function

\begin{verbatim}
def display_patient(patient):
    # your code here
\end{verbatim}

that prints the record in the form

\begin{verbatim}
ID: 1042
Name: Maria Garcia
Age: 37
\end{verbatim}

\begin{itemize}

\item In the function definition, what is the
\textbf{parameter}?

\item When you call \verb+display_patient(patient)+ what is the \textbf{argument}?

\end{itemize}

\end{enumerate}


\section{Returning a Result}

\begin{enumerate}

\item Consider the database

\begin{verbatim}
patients = [
 [1042, "Maria Garcia", 37],
 [1057, "Daniel Kim", 52],
 [1081, "Aisha Patel", 29]
]
\end{verbatim}

Write a function

\begin{verbatim}
find_patient(patients,patient_id)
\end{verbatim}

that searches through the database by patient ID.  If the ID is found, the function should \verb+return patient+.  If no record matches, use \verb+return None+.  A useful structure is

\begin{verbatim}
def find_patient(patients,patient_id):
    for patient in patients:
        if patient[0] == patient_id:
            return patient
    return None
\end{verbatim}

Why is \texttt{return} more useful here than simply printing the matching record inside the function?

\end{enumerate}


\section{Using a Module}

\begin{enumerate}

\item Place the two functions

\begin{verbatim}
display_patient()
find_patient()
\end{verbatim}

inside a file called

\begin{verbatim}
patient_tools.py
\end{verbatim}

In a second Python file or Google Colab notebook, import the
module:

\begin{verbatim}
import patient_tools
\end{verbatim}

Then search for patient ID \texttt{1057}:

\begin{verbatim}
result = patient_tools.find_patient(patients, 1057)
\end{verbatim}

Display the result using the function from the module:

\begin{verbatim}
patient_tools.display_patient(result)
\end{verbatim}

\begin{itemize}

\item Which file \textbf{defines} the functions?

\item Which file or notebook \textbf{calls} the functions?

\item What value is stored in \texttt{result}?

\end{itemize}

\end{enumerate}


\section{Think Ahead to Chapter 8}

\begin{enumerate}

\item Our current search requires an exact numeric patient ID.

Suppose instead that the user remembers only part of a patient's
name:

\begin{verbatim}
"Garcia"
\end{verbatim}

What new kind of operation would the program need in order to
locate the correct record?

\vspace{0.25cm}

\textbf{Takeaway:} Chapters 6--7 allow us to organize programs into
reusable functions and modules. Chapter 8 will allow us to make
the patient search more flexible by examining and matching strings.

\end{enumerate}

\end{document}
